\section{Related work}\label{sec:related}

\paragraph{Few-shot adaptation of vision-language models.}
Contrastive vision-language models such as CLIP \citep{radford2021clip} are adapted to new tasks with a few labelled examples in two broad ways.
Prompt learning replaces hand-written prompts with learned context vectors \citep{zhou2022coop}.
Adapter and residual methods instead keep the pretrained prediction and add a learned correction:
CLIP-Adapter blends adapted and original features through a residual connection \citep{gao2024clipadapter},
Tip-Adapter adds a cache built from the support set to CLIP's zero-shot logits \citep{zhang2022tipadapter},
and TaskRes learns a residual on top of the frozen text classifier \citep{yu2023taskres}.
All of these combine a pretrained and a learned component, and the weight given to each is usually a hyperparameter.
In biomedical imaging, BiomedCLIP \citep{zhang2023biomedclip} provides the pretrained model,
and BiomedCoOp \citep{koleilat2025biomedcoop} and vMFCoOp \citep{shao2026vmfcoop} adapt prompt learning to it;
BiomedCoOp's GPT-4 class descriptions are also the prompts PARA uses.

\paragraph{Setting the blending weight without validation data.}
With a few examples per class there is rarely data to spare for validation, so several methods set the balance between pretrained and learned components without it.
ProLIP fine-tunes only CLIP's image projection, regularized towards its pretrained value,
and in its validation-free variant scales the regularization weight as $1/N$ in the number of shots per class \citep{fahes2026prolip}.
HOSO-Adapter holds out one support example per class to learn the blending ratio of CLIP-Adapter \citep{vorster2026hoso}.
PARA \citep{kang2026para} computes its step in closed form from the geometry of the frozen support embeddings, with no training and no held-out data.
Our shot-only and HOSO-style baselines follow the first two ideas, applied to PARA's blend, so that the three ways of setting the step can be compared on the same trained models.

\paragraph{Geometry of contrastive embedding spaces.}
\citet{liang2022mind} showed that CLIP-style models place image and text embeddings in separate regions of the shared space (the modality gap),
with each modality confined to a narrow cone.
These two properties explain our first finding: PARA's discrepancy mostly measures the gap, which is the same for every class,
and its agreement mostly measures the cone, which is the same for any set of images.
We build on their analysis rather than claim the phenomenon as new.

\paragraph{Evaluation under support-set variation.}
Because few-shot results depend on the sampled support set, PARA evaluates over repeated support draws and reports the worst $20\%$ as well as the mean.
We adopt this protocol, share the draws across every method so that comparisons are paired,
and add balanced accuracy, tests per setting with multiple-comparison correction, and a split into design and held-out datasets.
